\documentclass[10pt,a4paper]{article}

\usepackage[utf8]{inputenc}
\usepackage[T1]{fontenc}
\usepackage[brazil]{babel}
\usepackage{amsmath, amssymb}
\usepackage[margin=1.2cm]{geometry}
\usepackage{enumitem}
\usepackage{booktabs}
\usepackage{parskip}

\setlength{\parskip}{5pt}
\setlength{\baselineskip}{1pt}

\newcommand{\R}{\mathrm{R}}
\newcommand{\A}{\mathrm{A}}
\newcommand{\MTTF}{\mathrm{MTTF}}
\newcommand{\MTBF}{\mathrm{MTBF}}
\newcommand{\MTTR}{\mathrm{MTTR}}

\setlist[enumerate,1]{leftmargin=*, label=\alph*), itemsep=0pt, topsep=1pt, parsep=0pt}
\setlist[itemize,1]{leftmargin=*, itemsep=0pt, topsep=1pt, parsep=0pt}

\begin{document}

\begin{center}
  {\large\bfseries Exercícios — Confiabilidade e Disponibilidade}\\[2pt]
  {\small Confiabilidade e Segurança de Software \textbar{} PUCRS}
\end{center}

\vspace{2pt}
\hrule
\vspace{4pt}

\textbf{Exercício 1: MTBF, MTTR e \(\A\)}

Um sistema operou por 30 dias. Houve 6 falhas registradas.\\
Tempos de reparo (h): 1,0 \ 0,5 \ 2,0 \ 1,5 \ 1,0 \ 0,5.
\begin{enumerate}
  \item Estime \(\MTBF\) em horas.
  \item Estime \(\MTTR\) em horas.
  \item Calcule \(\A\).
\end{enumerate}

\vspace{6pt}

\textbf{Exercício 2: Disponibilidade e SLA}

Uma plataforma tem \(\MTBF=200\,h\) e \(\MTTR=4\,h\).
\begin{enumerate}
  \item Calcule \(\A\).
  \item Estime o downtime por ano.
\end{enumerate}

\vspace{6pt}

\textbf{Exercício 3: \(\R(t)\) no modelo exponencial}

Assuma falhas exponenciais e \(\MTBF=500\,h\).
\begin{enumerate}
  \item Qual a probabilidade de operar sem falhar por 100\,h?
  \item E por 500\,h?
\end{enumerate}

\vspace{6pt}

\textbf{Exercício 4: Melhorar \(\MTBF\) ou \(\MTTR\)?}

Baseline: \(\MTBF=100\,h\), \(\MTTR=5\,h\).
\begin{itemize}
  \item Opção A: \(\MTBF \to 200\,h\) (mantendo \(\MTTR=5\,h\)).
  \item Opção B: \(\MTTR \to 1\,h\) (mantendo \(\MTBF=100\,h\)).
\end{itemize}
Qual opção melhora mais a disponibilidade?

\textbf{Exercício 5: Confiabilidade vs.\ disponibilidade}

Considere dois sistemas reparáveis:
\begin{itemize}
  \item Sistema A: \(\MTBF=400\,h\), \(\MTTR=8\,h\)
  \item Sistema B: \(\MTBF=250\,h\), \(\MTTR=2\,h\)
\end{itemize}
\begin{enumerate}
  \item Calcule a disponibilidade aproximada de cada sistema.
  \item Estime o downtime anual de cada um.
  \item Qual sistema é \textbf{mais confiável}?
  \item Qual sistema é \textbf{mais disponível}?
  \item Por que essas respostas não são necessariamente iguais?
\end{enumerate}

\vspace{6pt}

\textbf{Exercício 6: Confiabilidade em uma missão}

Um serviço deve operar por \(72\,h\) continuamente. Assuma modelo exponencial.
\begin{itemize}
  \item Sistema atual: \(\MTBF=600\,h\)
  \item Proposta de melhoria: \(\MTBF=900\,h\)
\end{itemize}
\begin{enumerate}
  \item Calcule \(\R(72)\) do sistema atual.
  \item Calcule \(\R(72)\) após a melhoria.
  \item Qual o ganho absoluto em pontos percentuais?
  \item A melhoria parece pequena ou relevante para uma missão crítica?
\end{enumerate}

\vspace{6pt}

\textbf{Exercício 7: Meta de SLA e requisito de \(\MTTR\)}

Uma plataforma possui \(\MTBF=300\,h\) e deseja cumprir \(\A \ge 99{,}5\%\).
\begin{enumerate}
  \item Qual deve ser o \(\MTTR\) máximo para atender a meta?
  \item Com esse \(\MTTR\), qual o downtime anual aproximado?
  \item Se na prática \(\MTTR=3\,h\), o SLA será atendido?
\end{enumerate}

\end{document}
